% !TeX program = lualatex
% Standalone research notebook for original catalogue problem 62.
% Compile twice: lualatex 62.tex
% Original entry retained; research subsection and [E] references added.
% No external bibliography, image, data, or style file is required.
\documentclass[11pt,a4paper]{article}
\usepackage[margin=24mm,headheight=14pt]{geometry}
\usepackage{fontspec}
\setmainfont{Latin Modern Roman}
\setsansfont{Latin Modern Sans}
\setmonofont{Latin Modern Mono}
\usepackage{amsmath,amssymb,mathtools,amsthm}
\usepackage{unicode-math}
\setmathfont{Latin Modern Math}
\usepackage{microtype}
\usepackage{enumitem,xurl,needspace}
\usepackage[dvipsnames]{xcolor}
\usepackage{hyperref}
\hypersetup{unicode=true,colorlinks=true,linkcolor=MidnightBlue,
 urlcolor=MidnightBlue,bookmarksnumbered=true,
 pdftitle={Research notebook - Problem 62},
 pdfauthor={Research attempt; not independently refereed}}
\usepackage{bookmark}
\usepackage{titlesec}
\titleformat{\section}{\Large\bfseries\raggedright}{\thesection}{0.7em}{}
\usepackage{fancyhdr}
\pagestyle{fancy}\fancyhf{}
\fancyhead[L]{\small\sffamily Research notebook}
\fancyhead[R]{\small\sffamily Problem 62}
\fancyfoot[L]{\scriptsize 27 September 2026}
\fancyfoot[C]{\thepage}
\fancyfoot[R]{\scriptsize Unrefereed research attempt}
\setlength{\parindent}{0pt}
\setlength{\parskip}{4pt plus 1pt}
\setlength{\emergencystretch}{3em}
\clubpenalty=10000
\widowpenalty=10000
\displaywidowpenalty=10000
\setcounter{secnumdepth}{1}
\setlist[itemize]{leftmargin=3em,itemsep=2pt,topsep=3pt,parsep=0pt,before=\small}
\allowdisplaybreaks
\newcommand{\N}{\mathbb{N}}
\newcommand{\Z}{\mathbb{Z}}
\newcommand{\Q}{\mathbb{Q}}
\newcommand{\R}{\mathbb{R}}
\newcommand{\C}{\mathbb{C}}
\newcommand{\F}{\mathbb{F}}
\newcommand{\A}{\mathbb{A}}
\newcommand{\PP}{\mathbb P}
\newcommand{\E}{\mathbb{E}}
\newcommand{\eps}{\varepsilon}
\newcommand{\dd}{\,\mathrm d}
\newcommand{\sref}[2]{\hyperlink{source:#1:#2}{[S#2]}}
\newcommand{\eref}[2]{\hyperlink{extra:#1:#2}{[E#2]}}
\newif\ifshowcatalogue
\showcataloguetrue
\newcommand{\cataloguescope}[1]{\ifshowcatalogue
 \paragraph{Exact scope in the supplied catalogue.}
 \begin{quote}\small #1\end{quote}\fi}
\newtheorem{notebooklemma}{Lemma}[section]
\newtheorem{notebookproposition}[notebooklemma]{Proposition}
\numberwithin{equation}{section}
\begin{document}
\setcounter{section}{61}
% ================================================================
% problemId: problem.top500-omission-w045049-dickson-linear-prime-values
% source releaseRank: 62; source releaseStatus: open_with_solved_subcases
% exactTarget SHA-256: 9109f0e7b5d0926524367a0a345e788e1fba6cfbdf9a820ea72426b2905d49b3
\clearpage
\section{\texorpdfstring{Dickson's Conjecture on Simultaneous Prime Values of Linear Forms}{Dickson's Conjecture on Simultaneous Prime Values of Linear Forms}}\label{62}
\subsection{Definitions and mathematical statement}
For $f_i(n)=a_in+b_i$ with integers $a_i\ge1$, define admissibility by the absence of a fixed prime divisor of $\prod_i f_i(n)$. Explicitly,
\[
 \forall p\in\mathcal P\ \exists n\in\Z:\quad p\nmid\prod_i(a_in+b_i).
\]
Dickson's assertion is $\forall X\ \exists n>X$ with all $a_in+b_i$ prime. The number of forms is finite and fixed for each assertion. Positivity for large $n$ follows from the positive slopes. The record does not require distinct forms or an asymptotic counting formula.
\par\smallskip\textit{Formulation and concept references:} \sref{62}{1}.
\cataloguescope{For every integer k>=1 and every family f\_i(n)=a\_i*n+b\_i (1<=i<=k) with integers a\_i>=1 and b\_i, if for each prime p some integer n has p not dividing the product of the f\_i(n), then infinitely many positive integers n make every f\_i(n) prime.}
\subsection{Short English statement}
Must every locally admissible finite family of positive-slope linear forms produce simultaneous primes arbitrarily often?
% BEGIN RESEARCH INSERTION
\subsection{Research attempt}

\paragraph{Current frontier.}
The formulation allows repeated forms and requires every form to be prime infinitely often \sref{62}{1}. Dirichlet's theorem settles one primitive positive-slope form. Maynard's multidimensional sieve produces many primes in sufficiently large admissible tuples, not simultaneous primality of every member of any prescribed tuple \eref{62}{1}. The other catalogue references are preserved; the listed proof-claim preprint is not used as a verified resolution.

\paragraph{Attack route.}
Better distribution in progressions can improve sieving, but does not by itself distinguish the required factor-count parity. Averaging over coefficients needs a separate uniformity argument to recover every fixed tuple. We pursue an exact parity expansion after sieving to numbers with at most two prime factors, and test whether lower-order sign moments could be enough.

\paragraph{Attempt: local preprocessing.}
Admissibility forces $\gcd(a_i,b_i)=1$ for each form. Delete repeated forms, which does not change the prime-producing set. Distinct remaining primitive coefficient pairs with positive slope cannot be proportional over $\Q$. Write $k$ for their number and
\[
 F(n)=\prod_{i=1}^k(a_in+b_i),\qquad
 \nu(p)=\#\{n\bmod p:F(n)=0\bmod p\}.
\]
Each form is nonzero modulo $p$ and has at most one root, so $\nu(p)\le k$. Once individual primitivity is known, admissibility therefore needs checking only for $p\le k$. For squarefree $d$, the Chinese remainder theorem gives
\begin{equation}\label{r62:crt}
 \#\{X<n\le2X:d\mid F(n)\}
 =\frac{\nu(d)}dX+O(\nu(d)),\qquad
 \nu(d)=\prod_{p\mid d}\nu(p).
\end{equation}
This follows by counting each residue class; it supplies no cancellation after a long sum over moduli.

\paragraph{An exact two-factor reduction.}
Fix $1/3<\theta<1/2$ and let $z=X^\theta$. For large $X$, all $f_i(n)$ in $X<n\le2X$ lie in $[cX,CX]$ for fixed $c,C>0$. Define
\[
 \mathcal R_X=\{n\in\Z:X<n\le2X,\ p\nmid F(n)
                       \text{ for all primes }p\le z\},\qquad R_X=|\mathcal R_X|.
\]
Every prime factor of each value on this set exceeds $z$. Three prime factors counted with multiplicity would give a value greater than $z^3>CX$. Hence
\[
 \Omega(f_i(n))\in\{1,2\},\qquad
 \symbf{1}_{\{f_i(n)\text{ prime}\}}=\frac{1-\lambda(f_i(n))}{2}.
\]
Prime squares have Liouville value $+1$ and are correctly excluded. Values at most one are absent for large $X$.

Let $P_X$ count simultaneous prime values in this interval and put
\[
 M_S(X)=\sum_{n\in\mathcal R_X}\prod_{i\in S}\lambda(f_i(n)),
 \quad S\subseteq[k],\qquad M_\varnothing=R_X.
\]
Every simultaneous-prime tuple is rough since $cX>z$. Expanding the exact indicator proves
\begin{equation}\label{r62:exact}
 P_X=2^{-k}\sum_{S\subseteq[k]}(-1)^{|S|}M_S(X).
\end{equation}

\paragraph{A bias-aware sufficient estimate.}
Sifting changes the ratio of primes to semiprimes; first moments need not vanish. Choose a fixed $a\in(0,1)$ and define
\[
 C_S(X)=\sum_{n\in\mathcal R_X}\prod_{i\in S}(\lambda(f_i(n))+a).
\]
The proposed sign mean is $-a$, not asserted here as an arithmetic fact. Algebra yields
\begin{equation}\label{r62:center}
 P_X=2^{-k}\left((1+a)^kR_X+
 \sum_{\varnothing\ne S\subseteq[k]}(-1)^{|S|}(1+a)^{k-|S|}C_S(X)\right).
\end{equation}
\begin{notebookproposition}\label{r62:criterion}
Suppose along an unbounded sequence of $X$ one has $R_X\to\infty$ and
\[
 |C_S(X)|\le\eta^{|S|}R_X\quad(\varnothing\ne S\subseteq[k])
\]
for fixed $a\in(0,1)$ and $0\le\eta<(2^{1/k}-1)(1+a)$. Then this tuple satisfies Dickson's conclusion.
\end{notebookproposition}
\begin{proof}
Taking absolute values after~\eqref{r62:center} gives
\[
 P_X\ge2^{-k}\bigl(2(1+a)^k-(1+a+\eta)^k\bigr)R_X.
\]
Its coefficient is strictly positive, so prime-producing integers exist at arbitrarily large scales.
\end{proof}
The required lower bound for the rough set and the full set of centered moment estimates are both unproved. Neither is inferred from~\eqref{r62:crt}.

\paragraph{Stress test: every proper-subset moment can look random.}
For $k\ge2$, put the uniform probability measure on sign vectors satisfying
\[
 \epsilon_1\cdots\epsilon_k=-(-1)^k.
\]
Every nonempty proper subset product has mean zero. Flip one coordinate in the subset and one outside it: this preserves the total product and reverses the subset product, pairing all vectors. Nevertheless the all-minus vector has probability zero. Perfect lower-order independence in this unbiased abstract model does not force the desired pattern. This is not an arithmetic counterexample or a refutation of the biased sufficient criterion; it identifies why its top-order information cannot simply be discarded.

The exponent $\theta>1/3$ is essential to the exact indicator. Without the exclusion of three factors, the same Liouville formula also counts products of three or five primes. Conversely, moving $\theta$ to $1/2$ cannot be inserted without rechecking the sieve estimates. The script tests~\eqref{r62:exact} at $\theta=2/5$, $X=1000,10000$, for $(n,n+2)$ and $(n,n+2,n+6)$, using exact factorization and all subset sums. It also checks the sign-model obstruction. No limiting conclusion is drawn.

\paragraph{Exact finite output.}
The computations produce the following counts; all prime counts agree with the full alternating moment expansion:
\[
\begin{array}{c|c|r|r}
 X&\text{offsets of the forms}&R_X&P_X\\ \hline
 1000&(0,2)&52&26\\
 1000&(0,2,6)&23&9\\
 10000&(0,2)&299&137\\
 10000&(0,2,6)&90&28
\end{array}
\]
These four finite identities neither prove the growth of $R_X$ nor certify the moment estimates along an unbounded sequence.

\paragraph{Outcome.}
\textbf{Outcome: Exploratory approach with unresolved gap.}

The attempt gives an exact rough-set parity expansion, a bias-aware sufficient condition, and a finite model showing why all proper-subset moments can miss the target event. The analytic lower bound and full moment estimates remain missing. This is a precise obstruction analysis, not a new theorem on prime values. Resolving all tuples would in particular settle twin primes and the linear part of the earlier Schinzel entry.

% END RESEARCH INSERTION
\subsection{Sources}
\begin{itemize}\raggedright
\item[\textnormal{[S1]}]\hypertarget{source:62:1}{} Shaohua Zhang, Notes on Dickson's Conjecture, arXiv:0906.3850v3 (2009).
\url{https://arxiv.org/html/0906.3850}.
{\small\textit{Catalogue formulation source.}}
\item[\textnormal{[S2]}]\hypertarget{source:62:2}{} An exposition on the supersimplicity of certain expansions of the additive group of the integers.
\url{https://arxiv.org/html/2502.03915}.
{\small\textit{Catalogue research/status reference; not a proof endorsement.}}
\item[\textnormal{[S3]}]\hypertarget{source:62:3}{} From Golomb to Schinzel.
\url{https://www.preprints.org/manuscript/202505.0651}.
{\small\textit{Catalogue research/status reference; not a proof endorsement.}}

% BEGIN ADDED RESEARCH REFERENCES
\item[\textnormal{[E1]}]\hypertarget{extra:62:1}{} James Maynard, \emph{Small gaps between primes}, Annals of Mathematics 181 (2015), 383--413; arXiv:1311.4600, main admissible-tuple results.
\url{https://arxiv.org/abs/1311.4600}.
{\small\textit{Several prime values in suitable large tuples, not all members of every fixed tuple. Consulted 27 September 2026.}}
% END ADDED RESEARCH REFERENCES
\end{itemize}

\end{document}
